\chapter{Wyniki eksperymentu}
\label{chap:Wynikieksperymentu}

Zostaną zaprezentowane wyniki związane z badaniami empirycznymi, które zostały omówione w poprzednim rozdziale. Będą one omówione względem pomiarów z oryginalnych prac naukowych, aby zweryfikować poprawność tez postawionych przez badaczy.

%\section{Empiryczna ocena efektywności algorytmów korekty}
%\label{sec:Ocena}
%Przedstawienie wyników własnego eksperymentu. Zestawienie skuteczności metod w oparciu o metryki statystyczne (RMSE, MAE) w różnych scenariuszach zakłóceń.

\section{Weryfikacja wybranych badań z literatury przedmiotu}
\label{sec:Weryfikacja}
W niniejszym rozdziale zaprezentowano wyniki empirycznej weryfikacji wybranych reguł metodycznych wyłonionych na etapie sformalizowanego przeglądu literatury (SLR). Głównym celem opisanych badań była próba wiernego odtworzenia opisanych w literaturze źródłowej protokołów eksperymentalnych oraz ocena faktycznej skuteczności proponowanych architektur w określonych warunkach.

\subsection{Wyniki badania Local Time Index}
\label{subsec:WynikiLTI}
W pierwszym eksperymencie oceniono wpływ metody uzupełniania braków na jakość prognozy modelu zadaniowego LSSVM.
Porównano klasyczne imputacje (średnia, hot-deck, AR) z metodą LTI rekomendowaną przez regułę SLR dla konfiguracji \textit{task model}~=~LSSVM oraz pasma braków $20$--$30\%$.
Jakość oceniano metryką NRMSE na sześciu szeregach (funkcje sine Tabela~\ref{tab:rule01-sine} i sinc Tabela~\ref{tab:rule01-sinc}, Mackey--Glass Tabela~\ref{tab:rule01-mackey},
obciążenie elektroenergetyczne Polski Tabela~\ref{tab:rule01-poland}, Sunspot Tabela~\ref{tab:rule01-sunspot} oraz piec gazowy Jenkins--Box Tabela~\ref{tab:rule01-gas}) przy losowych brakach MCAR na poziomach $R\in\{5,10,15,20,25,50\}\%$. Poniższe tabele przedstawiają uzyskane wartości NRMSE; przy interpretacji reguły szczególną uwagę zwraca się na wiersze $R=20\%$ oraz $R=25\%$, które mieszczą się w przedziale braków $20$--$30\%$ z reguły.

\begin{table}[htbp]
  \centering
  \caption{Porównanie NRMSE --- funkcji sine.}
  \label{tab:rule01-sine}
  \begin{tabular}{lcccc}
    \hline
    $R$ & mean & hot-deck & AR & LTI \\
    \hline
    5\% & 0.0623 & 0.0019 & 0.0011 & 0.0025 \\
    10\% & 0.0687 & 0.0029 & 0.0021 & 0.0032 \\
    15\% & 0.0897 & 0.0034 & 0.0053 & 0.0026 \\
    20\% & 0.1358 & 0.0038 & 0.0093 & 0.0033 \\
    25\% & 0.1378 & 0.0044 & 0.0201 & 0.0057 \\
    50\% & 0.3157 & 0.0091 & 0.0623 & 0.0112 \\
    \hline
  \end{tabular}
\end{table}

\begin{table}[htbp]
  \centering
  \caption{Porównanie NRMSE --- funkcji sinc.}
  \label{tab:rule01-sinc}
  \begin{tabular}{lcccc}
    \hline
    $R$ & mean & hot-deck & AR & LTI \\
    \hline
    5\% & 0.0363 & 0.0116 & 0.0114 & 0.0076 \\
    10\% & 0.0553 & 0.0118 & 0.0112 & 0.0079 \\
    15\% & 0.0587 & 0.0119 & 0.0117 & 0.0151 \\
    20\% & 0.0875 & 0.0121 & 0.0123 & 0.0098 \\
    25\% & 0.1025 & 0.0120 & 0.0159 & 0.0110 \\
    50\% & 0.1684 & 0.0126 & 0.0398 & 0.0151 \\
    \hline
  \end{tabular}
\end{table}

\begin{table}[htbp]
  \centering
  \caption{Porównanie NRMSE --- szeregu Mackey--Glass.}
  \label{tab:rule01-mackey}
  \begin{tabular}{lcccc}
    \hline
    $R$ & mean & hot-deck & AR & LTI \\
    \hline
    5\% & 0.0509 & 0.0466 & 0.0453 & 0.0309 \\
    10\% & 0.0615 & 0.0589 & 0.0633 & 0.0446 \\
    15\% & 0.0788 & 0.0795 & 0.0850 & 0.0551 \\
    20\% & 0.1058 & 0.1014 & 0.1078 & 0.0698 \\
    25\% & 0.1090 & 0.0994 & 0.1081 & 0.0860 \\
    50\% & 0.2041 & 0.1873 & 0.1993 & 0.2115 \\
    \hline
  \end{tabular}
\end{table}

\begin{table}[htbp]
  \centering
  \caption{Porównanie NRMSE --- obciążenia elektroenergetycznego (Polska).}
  \label{tab:rule01-poland}
  \begin{tabular}{lcccc}
    \hline
    $R$ & mean & hot-deck & AR & LTI \\
    \hline
    5\% & 0.0654 & 0.0617 & 0.0629 & 0.0684 \\
    10\% & 0.0738 & 0.0664 & 0.0721 & 0.0741 \\
    15\% & 0.0756 & 0.0669 & 0.0675 & 0.0863 \\
    20\% & 0.0923 & 0.0730 & 0.0820 & 0.0848 \\
    25\% & 0.1078 & 0.0712 & 0.0827 & 0.0957 \\
    50\% & 0.1411 & 0.1076 & 0.1105 & 0.1326 \\
    \hline
  \end{tabular}
\end{table}

\begin{table}[htbp]
  \centering
  \caption{Porównanie NRMSE --- szeregu Sunspot.}
  \label{tab:rule01-sunspot}
  \begin{tabular}{lcccc}
    \hline
    $R$ & mean & hot-deck & AR & LTI \\
    \hline
    5\% & 0.2914 & 0.2914 & 0.2917 & 0.2939 \\
    10\% & 0.2900 & 0.2911 & 0.2918 & 0.3047 \\
    15\% & 0.2946 & 0.2938 & 0.2956 & 0.3042 \\
    20\% & 0.2920 & 0.2910 & 0.2931 & 0.3134 \\
    25\% & 0.2939 & 0.2932 & 0.2952 & 0.2689 \\
    50\% & 0.2881 & 0.2917 & 0.2963 & 0.3321 \\
    \hline
  \end{tabular}
\end{table}

\begin{table}[htbp]
  \centering
  \caption{Porównanie NRMSE --- pieca gazowego Jenkins--Box.}
  \label{tab:rule01-gas}
  \begin{tabular}{lcccc}
    \hline
    $R$ & mean & hot-deck & AR & LTI \\
    \hline
    5\% & 0.0797 & 0.0620 & 0.0511 & 0.0512 \\
    10\% & 0.0954 & 0.0810 & 0.0591 & 0.0679 \\
    15\% & 0.0984 & 0.0694 & 0.0599 & 0.0644 \\
    20\% & 0.1151 & 0.0736 & 0.0717 & 0.0667 \\
    25\% & 0.1341 & 0.0759 & 0.0752 & 0.0748 \\
    50\% & 0.1418 & 0.0773 & 0.0934 & 0.1177 \\
    \hline
  \end{tabular}
\end{table}

\subsection{Wyniki eksperymentu RMD-FSE}
\label{subsec:WynikiRMDFSE}
W drugim eksperymencie oceniono metodę RMD-FSE na tle pojedynczych sieci FIR (ang.\ \textit{Finite Impulse Response}) uczonych po wcześniejszym uzupełnieniu braków (wygładzające funkcje sklejane, algorytm EM oraz jego wersja regularyzowana, w wariantach losowego oraz uśrednionego doboru oszacowań). Jakość predykcji oceniano za pomocą metryki NRMSE oraz wskaźnika $P_d$ [\%] wyznaczanego względem referencyjnej sieci \textbf{kompletnej FIR} (czyli sieci trenowanej na pełnym zbiorze uczącym, bez ubytków). Dla obu rozpatrywanych szeregów czasowych uzyskano stałe wartości NRMSE równe $0{,}0941$ (Sunspot) oraz $0{,}1468$ (RAMA PM2.5), przy wynikającej z definicji wartości $P_d = 0\%$. Wiersze odpowiadające kompletnej sieci FIR pominięto w tabelach, aby nie powielać wyników.


Dla każdej metody i poziomu braków $R$ wykonano $25$ niezależnych powtórzeń:
w każdym przebiegu losowo, zgodnie z mechanizmem MCAR, usuwano inną część obserwacji treningowych, po czym uczono model i mierzono błąd na teście. Raportowane NRMSE oraz Pd\,\%
są średnimi z tych $25$ realizacji. Poniższe zestawienia obejmują poziomy $R\in\{2{,}5;\,5;\,7{,}5;\,10;\,12{,}5;\,15\}\%$, Tabela~\ref{tab:rule02-sunspot} dla szeregu czasowego Sunspot, a Tabela~\ref{tab:rule02-rama} dla RAMA PM2.5, przy czym pasmo rozpatrywanej reguły z SLR $0$--$10\%$ odpowiada $R\in\{2{,}5;\,5;\,7{,}5\}$, co podsumowuje Tabela~\ref{tab:rule02-rmd-focus}.

\begin{table}[htbp]
  \centering
  \scriptsize
  \begin{minipage}[t]{0.48\textwidth}
    \centering
    \caption{Wyniki Pd\,\% i NRMSE --- Sunspot.}
    \label{tab:rule02-sunspot}
    \begin{tabular}{llrr}
      \hline
      $R$ & metoda & Pd\,\% & NRMSE \\
      \hline
      2.5\% & Spline & -6.22 & 0.0911  \\
     & EM-rand & -1.96 & 0.0931  \\
     & EM-avg & -1.85 & 0.0932 \\
     & RegEM-rand & -0.58 & 0.0938 \\
     & RegEM-avg & -2.03 & 0.0931 \\
     & RMD-FSE & -5.73 & 0.0913 \\
    \hline
    5\% & Spline & -7.03 & 0.0907 \\
     & EM-rand & -2.13 & 0.0930 \\
     & EM-avg & -0.83 & 0.0937 \\
     & RegEM-rand & -1.50 & 0.0933 \\
     & RegEM-avg & -3.04 & 0.0926 \\
     & RMD-FSE & -5.92 & 0.0912 \\
    \hline
    7.5\% & Spline & -7.44 & 0.0905 \\
     & EM-rand & -0.43 & 0.0938 \\
     & EM-avg & -2.21 & 0.0930 \\
     & RegEM-rand & 0.65 & 0.0944 \\
     & RegEM-avg & -3.49 & 0.0924 \\
     & RMD-FSE & -5.95 & 0.0912 \\
    \hline
    10\% & Spline & 2.81 & 0.0942 \\
     & EM-rand & 4.22 & 0.0955 \\
     & EM-avg & 3.80 & 0.0952 \\
     & RegEM-rand & 5.54 & 0.0963 \\
     & RegEM-avg & 3.42 & 0.0952 \\
     & RMD-FSE & 2.24 & 0.0943 \\
    \hline
    12.5\% & Spline & -4.66 & 0.0918 \\
     & EM-rand & 5.77 & 0.0965 \\
     & EM-avg & 2.31 & 0.0950 \\
     & RegEM-rand & 1.81 & 0.0948 \\
     & RegEM-avg & -0.69 & 0.0936 \\
     & RMD-FSE & -2.46 & 0.0928 \\
    \hline
    15\% & Spline & 3.71 & 0.0945 \\
     & EM-rand & 4.99 & 0.0959 \\
     & EM-avg & 6.72 & 0.0962 \\
     & RegEM-rand & 7.54 & 0.0967 \\
     & RegEM-avg & 7.80 & 0.0964 \\
     & RMD-FSE & 4.89 & 0.0953 \\
    \hline
    \end{tabular}
  \end{minipage}
  \hfill
  \begin{minipage}[t]{0.48\textwidth}
    \centering
    \caption{Wyniki Pd\,\% i NRMSE --- RAMA PM2.5.}
    \label{tab:rule02-rama}
    \begin{tabular}{llrr}
      \hline
      $R$ & metoda & Pd\,\% & NRMSE \\
      \hline
      2.5\% & Spline & 5.29 & 0.1505 \\
     & EM-rand & 0.03 & 0.1466 \\
     & EM-avg & 1.91 & 0.1480 \\
     & RegEM-rand & -0.83 & 0.1461 \\
     & RegEM-avg & 4.61 & 0.1500 \\
     & RMD-FSE & -15.97 & 0.1345 \\
    \hline
    5\% & Spline & 6.71 & 0.1514 \\
     & EM-rand & -0.91 & 0.1460 \\
     & EM-avg & 2.12 & 0.1481 \\
     & RegEM-rand & -0.74 & 0.1461 \\
     & RegEM-avg & 5.54 & 0.1506 \\
     & RMD-FSE & -14.76 & 0.1355 \\
    \hline
    7.5\% & Spline & 2.71 & 0.1486 \\
     & EM-rand & 0.35 & 0.1469 \\
     & EM-avg & 0.33 & 0.1468 \\
     & RegEM-rand & -0.19 & 0.1465 \\
     & RegEM-avg & 2.71 & 0.1486 \\
     & RMD-FSE & -15.73 & 0.1347\\
    \hline
    10\% & Spline & 2.92 & 0.1487 \\
     & EM-rand & 2.12 & 0.1481 \\
     & EM-avg & 0.33 & 0.1468 \\
     & RegEM-rand & 1.61 & 0.1478 \\
     & RegEM-avg & 0.96 & 0.1474 \\
     & RMD-FSE & -13.90 & 0.1361 \\
    \hline
    12.5\% & EM-rand & -1.15 & 0.1458 \\
     & EM-avg & 1.56 & 0.1476 \\
     & RegEM-rand & -0.82 & 0.1459 \\
     & RegEM-avg & 0.25 & 0.1468 \\
     & RMD-FSE & -15.97 & 0.1345 \\
    \hline
    15\% & Spline & 11.59 & 0.1539 \\
     & EM-rand & 3.12 & 0.1487 \\
     & EM-avg & 7.34 & 0.1517 \\
     & RegEM-rand & 2.16 & 0.1481 \\
     & RegEM-avg & 6.57 & 0.1511 \\
     & RMD-FSE & -6.23 & 0.1405\\
    \hline
    \end{tabular}
  \end{minipage}
\end{table}

\begin{table}[htbp]
  \centering
  \caption{Podsumowanie RMD-FSE w pasmie 0--10\%.}
  \label{tab:rule02-rmd-focus}
  \begin{tabular}{llcc}
    \hline
    zbiór & $R$ & Pd\,\% & NRMSE \\
    \hline
    Sunspot & 2.5\% & -5.73 & 0.0913 \\
    Sunspot & 5\% & -5.92 & 0.0912 \\
    Sunspot & 7.5\% & -5.95 & 0.0912 \\
    RAMA PM2.5 & 2.5\% & -15.97 & 0.1345 \\
    RAMA PM2.5 & 5\% & -14.76 & 0.1355 \\
    RAMA PM2.5 & 7.5\% & -15.73 & 0.1347 \\
    \hline
  \end{tabular}
\end{table}


\begin{table}[htbp]
  \centering
  \caption{Wyniki na ETTh1: Transformer (full) vs Transformer* (Decay Attention).}
  \label{tab:rule04-etth1}
  \begin{tabular}{lcccccc}
    \hline
    $\mathrm{pred\_len}$ & \multicolumn{2}{c}{MAE} & \multicolumn{2}{c}{MSE} & \multicolumn{2}{c}{$R^2$} \\
    \cline{2-7}
     & full & Decay & full & Decay & full & Decay \\
    \hline
    24 & 0.4952 & 0.4989 & 0.4690 & 0.4695 & 0.5335 & 0.5329 \\
    48 & 0.5516 & 0.5796 & 0.5530 & 0.6009 & 0.4500 & 0.4023 \\
    168 & 0.6920 & 0.6644 & 0.7659 & 0.7078 & 0.2371 & 0.2950 \\
    336 & 0.7493 & 0.7653 & 0.8786 & 0.8912 & 0.1171 & 0.1044 \\
    720 & 0.7266 & 0.7153 & 0.8385 & 0.8151 & 0.1443 & 0.1681 \\
    \hline
  \end{tabular}
\end{table}

\subsection{Wyniki badania Decay Attention}
\label{subsec:WynikiDA}
W trzecim eksperymencie porównano Transformer z pełną samo-uwagą (ang.\ \textit{full self-attention}) oraz wariant Transformer* z mechanizmem Decay Attention. Protokół artykułu źródłowego przewiduje cztery zbiory z rodziny ETT (ETTh1, ETTh2, ETTm1, ETTm2) oraz uśrednienie z pięciu niezależnych powtórzeń. Ze względu na koszt obliczeniowy treningu modeli Transformer weryfikację ograniczono do zbioru \textbf{ETTh1} oraz jednego powtórzenia ($\mathrm{itr}=1$), przy zachowaniu pozostałych elementów protokołu: długości okna $\mathrm{seq\_len}=96$, $\mathrm{label\_len}=48$ oraz horyzontów $\mathrm{pred\_len}\in\{24,48,168,336,720\}$. Tabela~\ref{tab:rule04-etth1} przedstawia uzyskane MAE, MSE oraz $R^2$ w porównaniu oryginalnego Transformera z wariantem z mechanizmem Decay Attention. 

\begin{table}[htbp]
  \centering
  \caption{Wyniki na ETTh1.}
  \label{tab:rule04-etth1}
  \begin{tabular}{lcccccc}
    \hline
    $\mathrm{pred\_len}$ & \multicolumn{2}{c}{MAE} & \multicolumn{2}{c}{MSE} & \multicolumn{2}{c}{$R^2$} \\
    \cline{2-7}
     & full & Decay & full & Decay & full & Decay \\
    \hline
    24 & 0.4952 & 0.4989 & 0.4690 & 0.4695 & 0.5335 & 0.5329 \\
    48 & 0.5516 & 0.5796 & 0.5530 & 0.6009 & 0.4500 & 0.4023 \\
    168 & 0.6920 & 0.6644 & 0.7659 & 0.7078 & 0.2371 & 0.2950 \\
    336 & 0.7493 & 0.7653 & 0.8786 & 0.8912 & 0.1171 & 0.1044 \\
    720 & 0.7266 & 0.7153 & 0.8385 & 0.8151 & 0.1443 & 0.1681 \\
    \hline
  \end{tabular}
\end{table}

\subsection{Wyniki eksperymentu TRF / VSF}
\label{subsec:WynikiTRF}
W czwartym eksperymencie oceniono rozwiązanie TRF (ang.\ \textit{Time-series Reconstruction Flows}) w zadaniu prognozowania przy niepełnym podzbiorze zmiennych (ang. \textit{Variable Subset Forecasting}, VSF). W tym eksperymencie duży koszt obliczeniowy pełnego przebiegu badania uniemożliwił jego wykonanie w pełnym zakresie. Trening modelu prognostycznego MTGNN (ang. \textit{Multivariate Time Series Graph Neural Network}), trening TRF z meta-uczeniem VAML (ang. \textit{Variable-Agnostic Meta-Learning}) oraz ewaluacja VSF (w szczególności FDW przy przeszukiwaniu banku wzorców treningowych) --- weryfikację ograniczono do zbioru \textbf{METR-LA}. Tabela~\ref{tab:rule05-metr-la} przedstawia MAE, RMSE oraz względny przyrost błędu względem wariantu Oracle.

\begin{table}[htbp]
  \centering
  \caption{Wyniki VSF na METR-LA.}
  \label{tab:rule05-metr-la}
  \begin{tabular}{lcccc}
    \hline
    ustawienie & MAE & RMSE & $\Delta$MAE\,\% & $\Delta$RMSE\,\% \\
    \hline
    Oracle & 4.7119 & 12.0927 & --- & --- \\
    Partial & 7.4910 & 14.7541 & 58.98 & 22.01 \\
    Reconstruct (TRF) & 10.4153 & 16.0610 & 121.04 & 32.82 \\
    FDW & 5.0447 & 12.5553 & 7.06 & 3.83 \\
    \hline
  \end{tabular}
\end{table}


\section{Wpływ jakości rekonstrukcji na realizację zadań prognostycznych}
\label{sec:Wpływ}
%Prezentacja wyników eksperymentu w kontekście modeli zadaniowych. Wykazanie bezpośredniej, ilościowej zależności między precyzją zastosowanej korekty a finalną celnością modeli predykcyjnych.